\chapter{Що обчислюють \nt{SERO}-нейрони}
\label{sec:serocomputer}
\begin{chaptergoals}
\item як ритмоводій із гальмівним рецептором дає НЕ та АБО-НЕ з окремих \nt{SERO}-нейронів;
\item чому викид в об'єм лишає лише кілька незалежних проводів, заданих $\ell\sim\sqrt{D\tau}$;
\item як концентрація згладжує імпульси в рівень, а від'ємний зворотний зв'язок міг би його тримати;
\item чому \nt{SERO} --- кандидат на роль горизонту планування і як горизонт задає терпіння.
\end{chaptergoals}

\section{Перший широкомовний канал}

Досі всі нейрони книги говорили адресною шиною: \nt{GLUT} і \nt{GABA} надсилали імпульси певним отримувачам, і ми з'ясували, що обчислює така мережа і якою мовою вона говорить. Тепер переходимо до другої шини з підрозділу~\ref{sec:buses} --- широкомовної. Її перший канал --- \nt{SERO}. Як і раніше, дозволимо собі лише те, що буває в живому організмі.

У цьому розділі з'являться три пристрої, якими згодом користуватимуться всі широкомовні канали: нейрон, який за рівного фонового притоку працює сам, доки його не зупинять; провід, який насправді є об'ємом тканини; і повільна концентрація замість окремих імпульсів. \nt{DOPA}, \nt{NORA} і \nt{ACET} у розділах~\ref{sec:dofa}--\ref{sec:acet} буде зібрано з тих самих деталей.

З погляду інженера \nt{SERO}-нейрон має чотири важливі відмінності.

\emph{Перша: знак вибирає отримувач.} Серотонін має рецептори обох знаків, і правило~\eqref{eq:dalesign} тут не діє: знак визначає не відправник, а рецептор отримувача (твердження~\ref{prop:receptor})~\cite{BarnesSharp1999}.

\emph{Друга: \nt{SERO}-нейрон працює сам, якщо є рівний фон.} У стані неспання він рівно, кілька разів на секунду, видає імпульси, і поштовх на кожен імпульс йому не потрібен: це ритмоводій. Але йому потрібен сталий фоновий приток. У зрізі мозку, де цього притоку немає, більшість таких клітин мовчить; рівний ритм повертається, якщо подати норадреналін через $\alpha_1$-рецептори, а їхня блокада його гасить~\cite{VandermaelenAghajanian1983}. В організмі цей фон надходить від \nt{NORA}. Для схеми це генератор, якому потрібне живлення: поки воно є, він працює сам, доки його не зупинить гальмування.

\emph{Третя: він повільний.} Майже всі рецептори серотоніну метаботропні: вони не відкривають канал самі, а запускають у клітині ланцюжок реакцій. Відповідь повільна: гальмівний струм через головний рецептор наростає і спадає в межах приблизно секунди~\cite{CourtneyFord2016} --- у десятки разів довше, ніж відповідь швидкого синапсу \nt{GLUT}.

\emph{Четверта: він викидає медіатор не в провід, а в об'єм.} В областях, куди йдуть виходи \nt{SERO}-нейронів, серотонін часто виходить не у вузьку щілину, а в міжклітинний простір і розпливається навколо~\cite{Bunin1998}. Отримувач відчуває концентрацію і не може знати, від якого відправника прийшла молекула.

За таких часів окремі імпульси вже неважливі, тому описуватимемо мережу частотами $r_j$ і концентраціями. Концентрація $c$ серотоніну в об'ємі, куди викидають медіатор нейрони $j$, зростає від викиду і спадає, бо медіатор відкачують назад:
\begin{empheq}[box=\keybox]{equation}
\tau_c\,\frac{\dd c}{\dd t} \;=\; -c \;+\; \kappa\sum_j r_j ,
\qquad
r_i \;=\; f_i\bigl(b_i + s_i\,g\,c\bigr), \quad s_i=\pm1 .
\label{eq:seroneuron}
\end{empheq}
Це ще один спосіб прочитати потік імпульсів, що доповнює твердження~\ref{prop:reader}: об'єм не бачить ні окремих імпульсів, ні їхнього порядку, а лише частоту, усереднену за час $\tau_c$. Перше рівняння --- те саме RC-коло, що й мембрана: $\tau_c$ --- час життя медіатора в об'ємі; прямо його не виміряно, і ми вважаємо його порядку секунди, $\kappa$ --- скільки медіатора дає один імпульс. Друге описує отримувача: $b_i$ --- його власний приток, у ритмоводія додатний (до нього входить і рівний фоновий вхід від \nt{NORA}); $g$ --- чутливість рецептора; знак $s_i$ вибирає отримувач; $f_i$ --- неспадна передавальна характеристика.

\section{НЕ одним нейроном}

Візьмімо ритмоводій із гальмівним рецептором, $s=-1$. Поки серотоніну навколо немає, він працює на власному притоці $b$. Коли серотонін з'являється, приток зменшується на $gc$, і якщо $gc>b$, нейрон замовкає. Це НЕ: вихід є, коли входу немає. На нього пішов один нейрон (рис.~\ref{fig:serogates}). Твердження~\ref{prop:monotone} тут незастосовне: воно вимагало додатних ваг.

Якщо на той самий ритмоводій діє серотонін із двох різних об'ємів, він мовчить, коли серотонін є хоча б в одному з них. Це АБО-НЕ, а з АБО-НЕ можна зібрати будь-яку логічну функцію. Двопровідний код \nt{GLUT}-мережі тут не потрібен: відсутність сигналу обчислюють нейрони, які працюють, доки їх не зупинять.

\begin{figure}[t]
\centering
\begin{tikzpicture}[>=Stealth,
  nrn/.style={circle,draw,very thick,fill=blue!6,minimum size=0.95cm,inner sep=0pt},
  pace/.style={nrn,double,double distance=1.2pt},
  inh/.style={-{Bar[width=7pt]},thick,red!70!black}]
\node[pace] (N1) at (0,0) {};
\draw[inh] (-1.6,0) node[left,black] {$c$} -- (N1);
\draw[->,very thick,red!70!black] (N1) -- (1.3,0) node[right,black] {$\bar c$};
\node[below=0.55cm] at (0,0) {\small НЕ};
\node[pace] (N2) at (5,0) {};
\draw[inh] (3.4,0.45) node[left,black] {$c_1$} -- (N2);
\draw[inh] (3.4,-0.45) node[left,black] {$c_2$} -- (N2);
\draw[->,very thick,red!70!black] (N2) -- (6.3,0) node[right,black] {$\overline{c_1\vee c_2}$};
\node[below=0.55cm] at (5,0) {\small АБО-НЕ};
\end{tikzpicture}
\caption{Логіка з \nt{SERO}-нейронів. Подвійний кружок --- ритмоводій: за рівного фонового притоку він працює сам, доки його не загальмують. Лінія з рискою --- гальмівний рецептор, $s=-1$; $c$, $c_1$, $c_2$ --- концентрації серотоніну в об'ємах навколо отримувача. Ліворуч --- НЕ, праворуч --- АБО-НЕ.}
\label{fig:serogates}
\end{figure}

\begin{keyblock}
\begin{proposition}[Повнота \nt{SERO}-логіки]
\label{prop:serocomplete}
Якщо в мережі є ритмоводії з гальмівними рецепторами і викиди в різні об'єми не змішуються, мережа~\eqref{eq:seroneuron} може обчислити будь-яку логічну функцію, і кожен біт передається одним проводом.
\end{proposition}
\end{keyblock}

Логічно \nt{SERO}-мережа сильніша за \nt{GLUT}, але умова «викиди в різні об'єми не змішуються» в мозку не виконується (підрозділ~\ref{sec:volume}).

\section{Від'ємний зворотний зв'язок}

Серотонінові нейрони мають гальмівні рецептори на собі самих~\cite{Sprouse1987}. Серотонін, який вони викинули, повертається до них же і пригальмовує їх. Для інженера це знайома схема: підсилювач із від'ємним зворотним зв'язком (рис.~\ref{fig:serofeedback}).

Що тут виміряно, а що --- модель. У самому ядрі шва, де розташовані \nt{SERO}-нейрони, серотонін гальмує сусідів не через спільний об'єм, а точка-точка, через синапси: переносник швидко прибирає його і не дає накопичитися поза щілиною. Поодинокий викид дає лише коротку паузу в розряді, а середня частота від цього не змінюється~\cite{CourtneyFord2016}. Тому петля нижче, замкнена через спільний об'єм, --- модель: ми рахуємо, що робив би такий зворотний зв'язок, якби він працював на рівні середніх частот.

\begin{figure}[t]
\centering
\begin{tikzpicture}[>=Stealth,
  nrn/.style={circle,draw,very thick,fill=blue!6,minimum size=0.95cm,inner sep=0pt},
  pace/.style={nrn,double,double distance=1.2pt},
  blk/.style={draw,very thick,rounded corners=2pt,fill=orange!10,minimum height=0.9cm,minimum width=2.2cm,align=center,font=\small},
  inh/.style={-{Bar[width=7pt]},thick,red!70!black}]
\node[pace] (N) at (0,0) {};
\node[blk] (V) at (3.6,0) {об'єм\\$\tau_c$};
\draw[->,very thick] (N) -- node[above] {\small $r$} (V);
\draw[->,very thick] (V) -- (6.0,0) node[right] {\small $c$ --- усім отримувачам};
\draw[inh] (V.south) -- ++(0,-0.9) -| node[pos=0.25,below] {\small $-g\,c$} (N.south);
\draw[->,thick,blue!60!black] (-1.7,0) node[left,black] {\small $b$} -- (N);
\end{tikzpicture}
\caption{\nt{SERO}-нейрон зі зворотним зв'язком через власний медіатор: концентрація $c$ іде до всіх отримувачів і гальмує сам нейрон. У моделі це регулятор рівня; у мозку така роль петлі --- гіпотеза.}
\label{fig:serofeedback}
\end{figure}

Порахуймо, що робить ця петля. Нехай передавальна характеристика лінійна, $f(u)=au$ при $u>0$. У рівновазі $c=\kappa r$ і $r=a(b-gc)$. Підставивши одне в інше, отримаємо
\begin{empheq}[box=\keybox]{equation}
r \;=\; \frac{a\,b}{1+G}, \qquad G \;=\; a\,g\,\kappa .
\label{eq:feedback}
\end{empheq}
Величина $G$ --- підсилення петлі. Це та сама формула, що й для будь-якого підсилювача з від'ємним зворотним зв'язком, і вона дає ті самі два виграші.

\emph{Стійкість до розкиду.} Нехай збудливість нейрона $a$ змінилася на частку $\varepsilon$. Без зворотного зв'язку на ту саму частку змінилася б і частота. Зі зворотним зв'язком, при $G\gg1$, частота $r\approx b/(g\kappa)$ майже не залежить від $a$: її зміна в $1+G$ разів менша. При $G=9$ розкид пригнічується вдесятеро.

\emph{Швидкодія.} Підставмо $r=a(b-gc)$ у перше рівняння~\eqref{eq:seroneuron}: отримаємо $\tau_c\,\dd c/\dd t=-(1+G)\,c+\kappa ab$. Концентрація наближається до свого рівня зі сталою часу $\tau_c/(1+G)$ --- в $1+G$ разів швидше, ніж без зворотного зв'язку.

Ось перше обчислення, яке могла б виконувати \nt{SERO}-система: тримати загальний рівень серотоніну в мозку на заданій позначці, як термостат тримає температуру. Це гіпотеза: у досліді автогальмування поки видно лише як короткі паузи, що не змінюють середньої частоти.

\section{Від імпульсів --- до рівня}

Друге обчислення сховане в першому рівнянні~\eqref{eq:seroneuron}. Нейрони видають окремі імпульси, а отримувач відчуває концентрацію, яка згладжує їх за час $\tau_c$. Якщо частота джерела змінюється, концентрація йде за нею із запізненням:
\begin{equation}
c(t) \;=\; \frac{\kappa}{\tau_c}\int_{-\infty}^{t} e^{-(t-t')/\tau_c}\,\sum_j r_j(t')\,\dd t' .
\label{eq:lowpass}
\end{equation}
Це ковзне середнє активності джерел за останні кілька $\tau_c$ --- фільтр нижніх частот (рис.~\ref{fig:lowpass}). Так найпростіший цифро-аналоговий перетворювач пропускає імпульси через RC-коло і перетворює їхню частоту на рівень. \nt{SERO}-система робить те саме із сотнями тисяч нейронів одночасно.

Ця величина водночас і пам'ять. Після того як джерела замовкли, концентрація тримається ще час порядку $\tau_c$. Це не біт, а слід, що плавно тане.

\begin{figure}[t]
\centering
\begin{tikzpicture}[>=Stealth,x=1.45cm,y=1cm]
\draw[gray!70!black] (0,3.2) -- (8.1,3.2);
\node[left,font=\small] at (-0.05,3.4) {імпульси};
\node[above,font=\scriptsize,gray!40!black] at (1,3.65) {$2$ на секунду};
\node[above,font=\scriptsize,gray!40!black] at (3.5,3.65) {$8$ на секунду};
\node[above,font=\scriptsize,gray!40!black] at (6.5,3.65) {$2$ на секунду};
\draw[->,gray!70!black] (0,0) -- (8.3,0) node[right,black] {\small $t$};
\draw[->,gray!70!black] (0,0) -- (0,2.75) node[left,black] {\small $c$};
\foreach \s in {0.136,0.529,0.760,1.223,1.714,1.748,1.755,2.663,2.701,2.734,3.414,3.493,3.719,3.800,3.928,3.948,4.074,4.327,4.420,4.589,4.728,4.736,4.914,5.025,5.205,5.220,6.224,6.544,7.178}{\draw[thick,blue!60!black] (\s,3.2) -- (\s,3.6);}
\foreach \a/\b/\v in {0.000/0.136/0.000,0.136/0.529/1.000,0.529/0.760/1.675,0.760/1.223/2.330,1.223/1.714/2.466,1.714/1.748/2.509,1.748/1.755/3.425,1.755/2.663/4.402,2.663/2.701/2.775,2.701/2.734/3.672,2.734/3.414/4.553,3.414/3.493/3.306,3.493/3.719/4.055,3.719/3.800/4.235,3.800/3.928/4.905,3.928/3.948/5.316,3.948/4.074/6.211,4.074/4.327/6.476,4.327/4.420/6.028,4.420/4.589/6.493,4.589/4.728/6.483,4.728/4.736/6.642,4.736/4.914/7.589,4.914/5.025/7.351,5.025/5.205/7.579,5.205/5.220/7.331,5.220/6.224/8.221,6.224/6.544/4.012,6.544/7.178/3.914,7.178/8.000/3.076}{\draw[very thick,orange!85!black] plot[domain=\a:\b,samples=10] (\x,{0.25*\v*exp(-(\x-\a))});}
\foreach \s/\p/\q in {0.136/0.000/1.000,0.529/0.675/1.675,0.760/1.330/2.330,1.223/1.466/2.466,1.714/1.509/2.509,1.748/2.425/3.425,1.755/3.402/4.402,2.663/1.775/2.775,2.701/2.672/3.672,2.734/3.553/4.553,3.414/2.306/3.306,3.493/3.055/4.055,3.719/3.235/4.235,3.800/3.905/4.905,3.928/4.316/5.316,3.948/5.211/6.211,4.074/5.476/6.476,4.327/5.028/6.028,4.420/5.493/6.493,4.589/5.483/6.483,4.728/5.642/6.642,4.736/6.589/7.589,4.914/6.351/7.351,5.025/6.579/7.579,5.205/6.331/7.331,5.220/7.221/8.221,6.224/3.012/4.012,6.544/2.914/3.914,7.178/2.076/3.076}{\draw[very thick,orange!85!black] (\s,0.25*\p) -- (\s,0.25*\q);}
\draw[thick,densely dashed,black] plot[domain=0:2,samples=30] (\x,{0.25*2*(1-exp(-\x))})
  plot[domain=2:5,samples=40] (\x,{0.25*(8-6.271*exp(-(\x-2)))})
  plot[domain=5:8,samples=40] (\x,{0.25*(2+5.688*exp(-(\x-5)))});
\foreach \x in {0,2,5,8}{\draw[gray!60!black] (\x,-0.05) -- (\x,0.05) node[below=3pt,font=\scriptsize] {\x\,с};}
\end{tikzpicture}
\caption{Від імпульсів --- до рівня. Угорі --- імпульси \nt{SERO}-нейронів: дві секунди рідко, три секунди часто, потім знову рідко. Унизу --- концентрація серотоніну~\eqref{eq:lowpass} при $\tau_c=1\,\text{с}$. Штрихова лінія --- те саме, якби імпульси йшли ідеально рівно.}
\label{fig:lowpass}
\end{figure}

\section{Провід --- це об'єм}
\label{sec:volume}

Повернімося до умови твердження~\ref{prop:serocomplete}. Серотонін, викинутий поблизу різними відправниками, складається в одну концентрацію.

\emph{Провід тут --- не волокно, а об'єм тканини.} Два сигнали лишаються різними, тільки якщо їх викидають в області, рознесені далі, ніж встигає розпливтися медіатор. За час $\tau$ молекула з коефіцієнтом дифузії $D$ відходить на
\begin{empheq}[box=\keybox]{equation}
\ell \;\sim\; \sqrt{D\,\tau}\, .
\label{eq:diffusionlength}
\end{empheq}
Звивистість міжклітинних щілин сповільнює дифузію невеликої молекули приблизно в $2.6$ раза~\cite{Sykova2008}, а коефіцієнт дифузії самого серотоніну в тканині не виміряно; за розміром молекули оцінимо його в кілька одиниць на $10^{-6}$~см$^2$/с. Якщо сигнал живе $\tau\sim1$~с (теж оцінка), то $\ell\sim\sqrt{3\cdot10^{-6}}$~см $\approx20$~мкм. Адреса в такій мережі --- це \emph{місце}: отримувач дізнається, від кого сигнал, лише з того, де він сам розташований.

Справжні \nt{SERO}-нейрони влаштовані так, що цих окремих адрес майже не лишається. Вихід кожного з них розкиданий по величезних ділянках мозку: гілки однієї клітини сягають багатьох далеких одна від одної областей~\cite{GagnonParent2014}. Місць викиду на клітину, за оцінкою, близько $10^5$ (таблиця~\ref{tab:types}); прямо їх не підраховано. «Проводи» різних \nt{SERO}-нейронів перекриваються і змішуються. Проте зовсім в один провід вони не зливаються: \nt{SERO}-нейрони поділяються на кілька підсистем, які надсилають вихід у різні області і по-різному відповідають на ту саму подію~\cite{Ren2018}. Але таких підсистем одиниці, а не сотні тисяч.

\begin{keyblock}
\begin{proposition}[Кількість незалежних проводів]
\label{prop:volumewires}
У мережі з об'ємною передачею кількість незалежних сигналів обмежена не кількістю нейронів, а кількістю неперетинних областей розміром $\ell\sim\sqrt{D\tau}$, у які вони викидають медіатор. Якщо вихід кожного нейрона розкиданий по великому об'єму, уся мережа передає лише кілька спільних змінних.
\end{proposition}
\end{keyblock}

Тому логічні схеми твердження~\ref{prop:serocomplete} в мозку нема з чого зібрати. А регуляторові~\eqref{eq:feedback} і фільтру~\eqref{eq:lowpass} окремі проводи не потрібні: їм досить однієї спільної величини. Фільтр прямо випливає з виміряного викиду в об'єм в областях призначення~\cite{Bunin1998}; регулятор рівня --- гіпотеза.

\section{Горизонт планування}
\label{sec:horizon}

Для чого може знадобитися одна спільна повільна величина? Добрий кандидат --- параметр, який має бути однаковим для всієї мережі і змінюватися не швидше, ніж за секунди чи хвилини. У розділі~\ref{sec:neurontypes} такий параметр уже з'являвся: коефіцієнт $\gamma$ у формулі помилки~\eqref{eq:rpe}, який показує, наскільки майбутня винагорода дешевша за негайну. У неперервному часі його місце займає \emph{горизонт} $\tau_\gamma$: винагорода $R$, на яку треба чекати час $D$, коштує зараз $R\,e^{-D/\tau_\gamma}$.

Горизонт вирішує, чи варто чекати. Нехай можна одразу взяти малу винагороду $r_0$ або дочекатися великої $R$. Чекати вигідно, доки
\begin{empheq}[box=\keybox]{equation}
D \;<\; \tau_\gamma\,\ln\frac{R}{r_0} .
\label{eq:patience}
\end{empheq}
При $\tau_\gamma=2\,\text{с}$ і втричі більшій відкладеній винагороді варто чекати до $2.2\,\text{с}$; якщо горизонт удвічі довший, --- до $4.4\,\text{с}$. Терпіння пропорційне горизонту.

Формула~\eqref{eq:patience} спирається на експоненційне знецінення. Виміряне знецінення на затримках у секунди ближче до гіперболи $R/(1+D/\tau_\gamma)$: так у мавп спадають і цінність відкладеної винагороди у виборі, і відповідь \nt{DOPA}-нейронів на її провісник~\cite{KobayashiSchultz2008}. Для гіперболи умова~\eqref{eq:patience} замінюється на $D<\tau_\gamma\,(R/r_0-1)$: залежність від відношення винагород уже не логарифмічна, а лінійна, і за тих самих чисел межа очікування зсувається майже вдвічі. Гіпербола виходить із тієї самої експоненти, якщо затримка випадкова (задача~\ref{ex:05:7}), а пропорційність терпіння масштабу часу зберігається.

За гіпотезою, горизонт і задає \nt{SERO}~\cite{Doya2002}. Для такої ролі в нього є все: один спільний рівень на всю мережу, що змінюється за секунди. Є й прямі досліди: якщо штучно ввімкнути серотонінові нейрони, тварина довше чекає на відкладену винагороду, перш ніж облишити очікування~\cite{Miyazaki2014}, і довше намагається здобути винагороду там, де вона стала рідкісною~\cite{Lottem2018}. Як саме концентрація серотоніну перетворюється на $\tau_\gamma$ усередині мережі, невідомо. У наступному розділі горизонт увійде в помилку, яку розсилає \nt{DOPA}.

\section{Підсумок}

\begin{table}[t]
\centering
\small
\begin{tabular}{>{\raggedright}p{3.6cm}>{\raggedright}p{4.3cm}>{\raggedright\arraybackslash}p{4.3cm}}
\toprule
& \nt{GLUT} & \nt{SERO} \\
\midrule
час реакції & мілісекунди & частки секунди --- секунда \\
знак дії & лише $+$ & $+$ і $-$, вибирає отримувач \\
без входу & мовчить & працює сам за рівного фонового притоку \\
адресація & точка --- точка & об'єм, адреса --- місце \\
НЕ & лише через датчики «вимкнення» & один нейрон \\
пам'ять & самопідтримна активність, гасне від утоми & концентрація, тане за час $\tau_c$ \\
головні обчислення & збіги, затримки, логіка, послідовності & згладжування; регулювання рівня і горизонт $\tau_\gamma$ (гіпотези) \\
\bottomrule
\end{tabular}
\caption{Що обчислюють мережі лише з \nt{GLUT}-нейронів і лише з \nt{SERO}-нейронів.}
\label{tab:glutvssero}
\end{table}

Як логічний елемент (таблиця~\ref{tab:glutvssero}) \nt{SERO}-нейрон багатший за \nt{GLUT}, але як система зв'язку він --- широкомовна розсилка: повільна і майже без адрес. Тому він не намагається бути процесором, а згладжує і розсилає кілька спільних величин, про які йшлося в розділі~\ref{sec:neurontypes}, --- за гіпотезою, зокрема й горизонт планування. Ритмоводій, об'єм і повільна концентрація трапляться нам ще тричі: у \nt{DOPA}, \nt{NORA} і \nt{ACET}.

\section{Задачі}

\begin{exercise}[\lvlA]
\label{ex:05:1}
\emph{Провід --- це об'єм.} Візьміть для серотоніну $D=3\cdot10^{-6}$~см$^2$/с (підрозділ~\ref{sec:volume}). Знайдіть за~\eqref{eq:diffusionlength} $\ell$ для сигналу, що живе $1\,\text{с}$, і час, за який медіатор розпливеться на $5$~см. Чим тоді забезпечено широкомовність \nt{SERO} --- дифузією чи розгалуженням проводу з $\sim10^5$ місць викиду?
\end{exercise}

\begin{exercise}[\lvlA]
\label{ex:05:2}
\emph{Регулятор рівня.} Покажіть, що в~\eqref{eq:seroneuron} при $f(u)=au$ відносна чутливість $S=(a/r)\,\partial r/\partial a$ дорівнює $1/(1+G)$ точно, а не лише при $G\gg1$. При $G=9$ збудливість $a$ зросла на $20\,\%$. На скільки відсотків зміниться $r$ без лінеаризації?
\end{exercise}

\begin{exercise}[\lvlB]
\label{ex:05:3}
\emph{Повільний рецептор у петлі.} Рецептор \nt{SERO} відповідає із затримкою: $r(t)=a\bigl(b-gc(t-T)\bigr)$, об'єм --- за~\eqref{eq:seroneuron}. Покажіть, що при $G>1$ петля стійка лише при $T<T^*=\tau_c\arccos(-1/G)/\sqrt{G^2-1}$. Знайдіть $T^*$ при $\tau_c=1\,\text{с}$ для $G=2$ і $G=9$. Яке пригнічення розкиду досяжне за затримки в сотні мілісекунд?
\end{exercise}

\begin{exercise}[\lvlB]
\label{ex:05:4}
\emph{Дробовий шум рівня.} Кожен імпульс додає до $c$ за~\eqref{eq:lowpass} експоненту з $\tau_c=1\,\text{с}$. Знайдіть коефіцієнт варіації $c$, якщо в один об'єм викидають пуассонівські імпульси всі $3\cdot10^5$ \nt{SERO}-нейронів по $2$ на секунду. Яке $\tau_c$ дало б $\mathrm{CV}=10\,\%$ одному нейрону з частотою $4$ імпульси на секунду?
\end{exercise}

\begin{exercise}[\lvlB]
\label{ex:05:5}
\emph{Моделювання: зворотний зв'язок проти шуму.} Змоделюйте $N=50$ пуассонівських ритмоводіїв із частотами $r_i=a_i[b-gc]_+$, $a_i$ рівномірні в $[2.8,\,5.2]$, і об'єм~\eqref{eq:seroneuron} з $\kappa=1/N$, $\tau_c=1\,\text{с}$. У скільки разів зворотний зв'язок ($b=1$, $g=2.25$) знижує CV рівня $c$ порівняно з $b=0.1$, $g=0$? Чи вирівнює він частоти нейронів?
\end{exercise}

\begin{exercise}[\lvlB]
\label{ex:05:6}
\emph{Проєктування: XOR на \nt{SERO}-логіці.} Вентиль --- ритмоводій, що видає $r_0$, якщо $b-g\sum_kc_k>0$, і $0$ інакше, у свій об'єм $\tau_c\,\dd c/\dd t=-c+\kappa r$; він обчислює АБО-НЕ. Зберіть XOR із п'яти таких вентилів і знайдіть час його встановлення при $\tau_c=1\,\text{с}$ і $G'=g\kappa r_0/b=2$.
\end{exercise}

\begin{exercise}[\lvlB]
\label{ex:05:7}
\emph{Терпіння за невідомої затримки.} (а)~Тварина погоджується чекати втричі більшу винагороду, якщо очікування не довше за $6\,\text{с}$. Який горизонт $\tau_\gamma$ це означає? (б)~Нехай затримка $D$ розподілена експоненційно із середнім $\bar D$. Покажіть, що чекати вигідно при $\bar D<\tau_\gamma(R/r_0-1)$, і порівняйте з фіксованою затримкою при $\tau_\gamma=2\,\text{с}$, $R=3r_0$.
\end{exercise}

\begin{exercise}[\lvlC]
\label{ex:05:8}
\emph{Як концентрація задає горизонт.} \nt{DOPA}-нейрон отримує $V$ напряму з вагою $k_1$ і через \nt{GABA}-посередника, що згладжує з $\tau_d=0.1\,\text{с}$, з вагою $-k_2$; для повільних $V$ сума дорівнює $(k_1-k_2)V+k_2\tau_d\dot V$. Доберіть $k_1$, $k_2$ під~\eqref{eq:ctd} з $\tau_\gamma=2\,\text{с}$. \nt{SERO} множить $k_1$ на $m$: наскільки треба змінити $m$, щоб подвоїти горизонт?
\end{exercise}
